\section{Asymptotic expansions of real radiation}
Throughout this section $X=s_{\mathrm{phys}}/m^2-1>0$ and $s=1+X$.
The labels $UV$ and $UW$ denote the angular denominator pairs
defined in the real-emission calculation.

\subsection{Threshold terms of the UV sector}
The analytic inner integral permits a uniform threshold expansion on a
fixed interval.  Set $v=1+Xw$, with $0\leq w\leq1$ and
$\dd v=X\,\dd w$.  Then
\begin{align}
 \ln v&=\ln(1+Xw),\\
 \ln(1-k)&=\ln[1+Xw(1-w)]-\ln(1+Xw),\qquad
 k=\frac{Xw^2}{1+Xw}.
\end{align}
Both logarithms have uniformly convergent Taylor series for $0<X<1$.
The rational denominators left after the soft subtraction contain
$1+Xw$ and $1+Xw(1-w)$, which are bounded away from zero.
The apparent powers of $w^{-1}$ cancel when the Laurent coefficients
are combined.  This cancellation is performed before integrating the
coefficient polynomials.  Keeping the Jacobian is essential to obtain
the correct leading power.

Let $F_{UV}$ denote the finite Laurent coefficient before multiplication
by the real cross-section prefactor $1/[8s X]$.  Direct expansion and
polynomial integration give
\begin{align}
 F_{UV}(X)={}&\frac{128}{X}+\frac{608}{3}-\frac{64}{3}\ln X\\
 &+X\left(\frac{1568}{15}+\frac{64}{15}\ln X\right)
 +X^2\left(\frac{380}{9}+\frac{56}{15}\ln X\right)\\
 &+X^3\left(\frac{1364}{315}-\frac{400}{63}\ln X\right)
                    +O\!\left(X^4[1+|\ln X|]\right).
\end{align}
The inverse power belongs to the sector decomposition, not to the
physical threshold cross section.  It must be added to the other real
sectors before their common prefactor is expanded.  The stated remainder
becomes $O(X^3[1+|\ln X|])$ after that prefactor.
The exact coefficient calculation is recorded in
\nolinkurl{soft_asymptotic_UV_low.json}.

\subsection{Threshold terms of the UW sector}
For the $UW$ sector the existing rectangle $0<a<X$, $0<t<1$ has already
been integrated in $a$.  Its logarithms are expanded uniformly in $X$.
The difference of dilogarithms in its remaining kernel can be expanded
without separately expanding either divergent argument at $t=0$:
\begin{align}
 d_3(X,t)&=\operatorname{Li}_2\!\left(-\frac{(1-t)(1+Xt)}t\right)
                  -\operatorname{Li}_2\!\left(-\frac{1-t}t\right),\\
 \frac{\partial d_3}{\partial X}
 &=-\frac{t}{1+Xt}\left[-\ln t+\ln(1+Xt(1-t))\right],
 \qquad d_3(0,t)=0.
\end{align}
Both dilogarithm arguments are negative, which fixes their real branches.
The Taylor series of this derivative is uniform apart from the integrable
factor $t\ln t$.  The analogous derivative of
$\operatorname{Li}_2(-1-Xt)+\pi^2/12$ is
$-t\ln(2+Xt)/(1+Xt)$, also with a zero boundary value at $X=0$.
The apparent endpoint poles cancel when the complete coefficient at each
power of $X$ is collected, before its final $t$ integration.

To control the remainder, first apply the endpoint Taylor subtractions
of Section~\ref{sec:comp-UW-simplified} to the eight-function kernel
$\Phi_{UW}(X,t)$. Separate its explicit dependence on $\ln X$, which
is at most linear. The accompanying rational coefficients have poles
of order at most four at $X=0$. Define
\begin{equation}
 X^4\Phi_{UW}(X,t)=\Psi_0(X,t)+\ln X\,\Psi_1(X,t).
\end{equation}
For any fixed $0<X_0<1$, both functions $\Psi_j$ extend analytically to
a neighborhood of $[0,X_0]\times[0,1]$. Indeed, the nonendpoint
denominators $1+Xt$, $2+Xt$ and $1+Xt(1-t)$ are bounded away from zero
there. Each endpoint Taylor subtraction is an analytic divided
difference. For the function $J$, its short dilogarithm expression
establishes joint analyticity near $t=0$, and its original
$d_3-\ell\ln t+\ell p$ expression does so near $t=1$.
The eighth $X$ derivatives of the $\Psi_j$ are consequently bounded
on this compact rectangle. Taylor expansion through $X^7$ has a
remainder bounded uniformly in $t$ by $C X^8[1+|\ln X|]$, with $C$
independent of $X$ and $t$. Dividing by $X^4$ and integrating over the
fixed unit interval gives the remainder below. The coefficient
integrals give, for one $UW$ sector,
\begin{align}
 F_{UW}(X)={}&X\left(\frac{16}{15}-\frac{64}{5}\ln X\right)
 +X^2\left(-\frac43-\frac{16}{5}\ln X\right)\\
 &+X^3\left(\frac{788}{315}-\frac{236}{105}\ln X\right)
                  +O\!\left(X^4[1+|\ln X|]\right).
\end{align}
The normalization is again before $1/[8s X]$.  The complete real answer
contains $2F_{UW}$ and the separate rational Bose-restoration term.
The exact coefficients are saved in \nolinkurl{soft_asymptotic_UW_low.json}.

\subsection{The combined real contribution at low energy}
Put $X=s-1>0$. The finite real coefficient $F_{\mathrm R}^{\eta}$
in this paragraph is the constant term obtained with the symmetric
regulator $(X_1X_2)^\eta$, before matching to the photon-mass
calculation. With $P=p+k$, these are exactly the dimensionless
Dalitz variables
\begin{equation}
 X_1=x=\frac{2P\cdot k_1}{m^2},\qquad
 X_2=y=\frac{2P\cdot k_2}{m^2}.
\end{equation}
In the center-of-mass frame, $X_i=2\sqrt{s}\,\omega_i/m$.
Thus $(X_1X_2)^\eta=(xy)^\eta$ uses the same finite-part
normalization as the earlier integration. The subscripts distinguish
the two outgoing energies from the incident-energy parameter $X$.

For the regular rational, energy-logarithm, $W$, and Bose terms,
use reciprocal symmetry to replace the last energy interval by
$1<u<s$, with an overall factor $1/2$. Set
\begin{equation}
 u=1+Xw,\qquad 0<w<1,\qquad
 \ell=\ln\frac{w}{1-w}.
\end{equation}
The logarithms in the exact kernels become
\begin{align}
 L_1&=2\ln(1+Xw)-\ln(1+X),\\
 L_2&=\ell+\ln(1+Xw),\\
 L_3&=\ell+2\ln[2+X(1+w)]-\ln(1+Xw)-2\ln(2+Xw).
\end{align}
Their expansions now have a fixed integration interval. The
removable Gram factors are combined before expanding. For example,
$\ln(a/b)/(a-b)$ is analytic at $a=b>0$. Such divided differences
remove the apparent interior pole near $w=1/2$.
After the four sectors are added, every coefficient through the
required order is a polynomial in $w$ multiplying $\ell$, plus a
rational function whose only poles are at the fixed endpoints.
This cancellation is an exact rational identity.

Subtract the simple and double endpoint terms before interchanging
the expansion and the integral. The remaining rational
denominators are bounded away from zero for sufficiently small
positive $X$ on the compact $w$ interval. Their subtracted
remainders are bounded by a constant times
$1+|\ln w|+|\ln(1-w)|$. This integrable bound justifies the
coefficientwise expansion. The endpoint counterterms themselves
are integrated exactly in the same finite-part convention.
The needed polynomial moment is
\begin{equation}
 \int_0^1w^n\ell\,\dd w=\frac{H_n}{n+1},\qquad
 H_n=\sum_{j=1}^{n}\frac1j,\qquad H_0=0.
\end{equation}
The finite parts of $w^{-1}$ and $(w-1)^{-1}$ vanish. Those of
$w^{-2}$ and $(w-1)^{-2}$ both equal $-1$.

Let $\Phi=\operatorname{FP}\int_0^1X f(1+Xw)\,\dd w$
for the sum of these four sectors. Let $A$ be its upper simple-pole
coefficient and $D_w$ its upper double-pole coefficient. Their
exact expansions are
\begin{align}
 \Phi={}&-\frac{192}{X}-\frac{1024}{3}-\frac{1792}{9}X
             -\frac{532}{9}X^2-\frac{7052}{135}X^3+O(X^4),\\
 A={}&\frac{32}{3}+\frac{32}{3}X+\frac43X^2
                            +\frac{188}{15}X^3+O(X^4),\\
 D_w={}&\frac{32}{X}+\frac{224}{3}+\frac{544}{9}X
                     +\frac{208}{9}X^2+\frac{172}{45}X^3+O(X^4).
\end{align}
The double pole belongs to the Bose-restoring contribution. Its
finite boundary correction must be included even though the
complete real integral eventually has only a logarithmic soft
divergence.

Denote the sum of the already derived $z$-cut to $\eta$-scheme
matching constants by $M_\eta$. Expanding their exact logarithms
and dilogarithms gives
\begin{equation}
 M_\eta=\frac{80}{3}+\frac{112}{3}X+\frac{62}{9}X^2
                              +\frac{352}{45}X^3+O(X^4\ln X).
\end{equation}
The unnormalized four-sector finite part is therefore
\begin{equation}
 F_4=\frac{\Phi}{2}+2A\ln X-\frac A2\ln(1+X)
                     -\frac{D_w(X+2)}{2(X+1)}+M_\eta.
\end{equation}
The logarithmic terms follow from the common $z$ boundary. The
term containing $D_w$ follows from the second-order displacement
of that boundary. Substitution yields
\begin{align}
 F_4={}&-\frac{128}{X}-\frac{608}{3}+\frac{64}{3}\ln X\\
 &+X\left(\frac{64}{3}\ln X-\frac{320}{3}\right)
 +X^2\left(\frac83\ln X-\frac{356}{9}\right)\\
 &+X^3\left(\frac{376}{15}\ln X-\frac{2596}{135}\right)
                                     +O(X^4\ln X).
\end{align}
Combining this with the independently expanded $UV$ and twice
the $UW$ contributions cancels every $X^{-1}$, constant, $X$,
and $X^2$ term, including their logarithms. The exact remaining
unnormalized expression is
\begin{equation}
 F_4+F_{UV}+2F_{UW}
 =X^3\left(\frac{128}{9}\ln X-\frac{1336}{135}\right)
                                      +O(X^4\ln X).
\end{equation}
Restoring the common phase-space factor $1/[8X(1+X)]$ gives
\begin{equation}
 \boxed{F_{\mathrm R}^{\eta}
 =X^2\left(\frac{16}{9}\ln X-\frac{167}{135}\right)
                                      +O(X^3\ln X).}
\end{equation}
The remainder follows from the uniform subtracted expansion and
the single explicit logarithm in the regulator matching. It is
not inferred by fitting numerical values. In the center-of-mass
local-subtraction convention, using
$A_{\mathrm{phys}}=8X^2/9+O(X^3)$ and
$\mathcal C_{\mathrm R}=F_{\mathrm R}^{\eta}
 -A_{\mathrm{phys}}\ln[2\sqrt{1+X}\,X]$, the same result is
\begin{equation}
 \mathcal C_{\mathrm R}
 =X^2\left[\frac89\ln\frac{X}{2}-\frac{167}{135}\right]
                                      +O(X^3\ln X).
\end{equation}

\subsection{The high-energy UV contribution}
Put $L=\ln X$. The constant-letter representation of the $UV$ finite coefficient gives
an independent route to its large-$X$ expansion.  In that representation,
$q=(\sqrt{X+4}-\sqrt{X})/(\sqrt{X+4}+\sqrt{X})$ tends to zero and
\begin{equation}
 q=\frac1X-\frac2{X^2}+\frac5{X^3}-\frac{14}{X^4}
                       +\frac{42}{X^5}+O(X^{-6}).
\end{equation}
The GPL differential equation expands each word in powers of $q$ and
$\ln q$.  The only singular kernel at this endpoint is $\dd q/q$.
All other kernels have convergent Taylor series about zero.  Words of
weight at most three therefore have remainders bounded by powers of
$q$ times at most $|\ln q|^3$.  The rational coefficients have poles
of order at most three, so the word expansion is carried through $q^4$
before extracting the physical $1/X$ term.

The integration constants in this expansion are the finite parts
$\operatorname{FP}G_1(w;0)$, where the subscript denotes basepoint one
and the same real left tangent as above.  They are evaluated analytically,
not determined by fitting.  Let
$a_\pm=(1\pm\ii\sqrt3)/2$ and $D(t)=1-t+t^2$.
The weight-one sum is
\begin{equation}
 G_1(a_+;t)+G_1(a_-;t)=\ln D(t).
\end{equation}
Shuffle identities remove most higher-weight endpoint constants.
The remaining weight-two sum is
\begin{equation}
 \sum_{a=a_\pm}G_1(1,a;0)
 =\int_0^1\frac{\ln D(t)}{1-t}\dd t
 =-\frac{\pi^2}{18}.
\end{equation}
The substitution $t\mapsto1-t$ reduces the integral to
$-\operatorname{Li}_2(a_+)-\operatorname{Li}_2(a_-)$.
The standard real boundary values follow directly by integrating the
Fourier series for $\ln(1-e^{\ii\theta})$, or by the dilogarithm
reflection identity on the conjugate pair.

Define the only remaining weight-three combination by
\begin{align}
 U&=\sum_{a,b=a_\pm}G_1(0,a,b;0)
                 =-\frac12\int_0^1\frac{\ln^2D(t)}t\dd t,\\
 V&=\sum_{a=a_\pm}G_1(-1,0,a;0)
                 =\int_0^1\frac{\ln D(t)\ln(1+t)}t\dd t.
\end{align}
Both are convergent real integrals.  The algebraic identity
$(1+t)D(t)=1+t^3$ gives
\begin{equation}
 V-U=\frac12\int_0^1
 \frac{\ln^2(1+t^3)-\ln^2(1+t)}t\dd t
 =-\frac13\int_0^1\frac{\ln^2(1+t)}t\dd t
 =-\frac{\zeta(3)}{12}.
\end{equation}
For the last elementary polylogarithmic moment, set $w=1/(1+t)$.
A primitive of $\ln^2w/(1-w)$ is
$-\ln^2w\ln(1-w)-2\ln w\operatorname{Li}_2(w)+2\operatorname{Li}_3(w)$.
Its values at $w=1/2,1$, together with
\begin{align}
 \operatorname{Li}_2(1/2)&=\frac{\pi^2}{12}-\frac12\ln^22,\\
 \operatorname{Li}_3(1/2)&=\frac78\zeta(3)-\frac{\pi^2}{12}\ln2
                                      +\frac16\ln^32,
\end{align}
prove that the moment is $\zeta(3)/4$.  Thus every boundary constant
needed here is fixed analytically.

Writing $\widehat F_{UV}=F_{UV}/[8s X]$, the result is
\begin{align}
 \widehat F_{UV}={}&X^3\left(\frac L{12}-\frac5{144}\right)
 +X^2\left(\frac L3+\frac1{18}\right)\\
 &+X\left(\frac5{24}L^2-\frac{49}{24}L+\frac{65}{72}\right)
 -\frac76L^2+\frac{11}{3}L+\frac{47}{36}+\frac{\pi^2}{6}\\
 &+\frac1X\left[7L^3-\frac{73}{6}L^2
 +\left(\frac{133}{8}-\frac{\pi^2}{3}\right)L
 +\frac{2335}{144}-\frac{\pi^2}{6}-4\zeta(3)\right]
 +O\!\left(\frac{L^3}{X^2}\right).
\end{align}
The growing terms are artifacts of the sector decomposition and must
cancel in the complete real contribution.  They are included explicitly
so that this cancellation can be checked rather than assumed.
The recurrence and the exact endpoint-constant reductions are implemented
in \nolinkurl{compton_soft_asymptotic_uv_high.py}, with the reduced result
saved in \nolinkurl{soft_asymptotic_UV_high.json}.

\section{The total cross section at low energy}
\label{sec:comp-independent-total-low}
The preceding expansions determine the first nonzero radiative
correction to the total cross section. Here
$X=s_{\mathrm{phys}}/m^2-1=2\omega/m$, with $\omega$ the incident
photon energy in the electron rest frame, and
$\ell_\lambda=\ln(\lambda_{\mathrm{phys}}/m)$.
Before taking the low-energy limit, cancel the exact photon-mass
logarithm in \eqref{eq:comp-assembly-common-regulator} by writing
\begin{equation}
 C_1=\bigl[C_V-A\ell_\lambda\bigr]+F_{\mathrm R}^{\eta}
       -A\ln\bigl[X\sqrt{1+X}\bigr]-\frac{\mathcal C_K}{1+X}.
\end{equation}
Each term in this expression is independent of the photon mass.
Consequently the following expansion does not interchange a
fixed-regulator threshold expansion with the infrared limit:
\begin{align}
 C_V-A\ell_\lambda
 &=X^2\left[-\frac83\ln X+\frac{22}{9}\right]
       +O\!\left(X^3|\ln X|\right),\\
 F_{\mathrm R}^{\eta}
 &=X^2\left[\frac{16}{9}\ln X-\frac{167}{135}\right]
       +O\!\left(X^3|\ln X|\right),\\
 -A\ln\bigl[X\sqrt{1+X}\bigr]-\frac{\mathcal C_K}{1+X}
 &=X^2\left[-\frac89\ln X-\frac{20}{27}\right]
       +O\!\left(X^3|\ln X|\right).
\end{align}
The logarithmic coefficients sum to $-16/9$, and the constants sum
to $22/9-167/135-20/27=7/15$. Thus
\begin{equation}
 \boxed{\sigma_1=\frac{\alpha_{\mathrm T}^3}{m^2}
 \left\{X^2\left[-\frac{16}{9}\ln X+\frac7{15}\right]
       +O\!\left(X^3|\ln X|\right)\right\},\qquad X\to0^+.}
 \label{eq:comp-independent-total-low}
\end{equation}
This is the order-$\alpha_{\mathrm T}^3$ contribution. The NLO total
is $\sigma_0+\sigma_1$, with the Born expansion
\begin{equation}
 \sigma_0=\frac{\pi\alpha_{\mathrm T}^2}{m^2}
 \left[\frac83-\frac83X+\frac{52}{15}X^2
                   -\frac{133}{30}X^3+O(X^4)\right].
\end{equation}
The remainder follows by adding the remainders of the independently
expanded, regulator-independent terms. No numerical fit or published
total is used to fix either coefficient. Since $X^2\ln X\to0$, this
order of the radiative correction vanishes in the Thomson limit.
\subsection{Expansion of the collected analytic formula}
\label{sec:comp-total-energy-low}
The complete product form \eqref{eq:comp-total-energy-product} gives
two further orders beyond the leading result above. The external energy variable remains $X$.
Direct expansion gives
\begin{equation}
\begin{aligned}
 \sigma_1=\frac{\alpha_{\mathrm T}^3}{m^2}\biggl\{
 &X^2\left[-\frac{16}{9}\ln X+\frac7{15}\right]
 +X^3\left[\frac43\ln X-\frac{113}{90}\right]\\
 &+X^4\left[-\frac{25}{9}\ln X+\frac{203}{125}\right]
 +O\!\left(X^5|\ln X|\right)\biggr\},\qquad X\to0^+.
\end{aligned}
\label{eq:comp-collected-total-low-four}
\end{equation}
No individual virtual or real integral is expanded again here. The
starting point is the complete infrared-finite analytic formula.

The expansion can be performed recursively from the defining GPL
kernels. For a nonzero rational letter $a$, and for the root letter,
\begin{align}
 \frac1{X-a}&=-\sum_{j=0}^{\infty}\frac{X^j}{a^{j+1}},
                      \qquad |X|<|a|,\\
 \frac1{\sqrt{X(X+4)}}
 &=\frac1{2\sqrt X}\sum_{j=0}^{\infty}
       \binom{-1/2}{j}\frac{X^j}{4^j},\qquad |X|<4.
\end{align}
For the zero letter the kernel is $1/X$. Starting with the empty
word $G(\varnothing;X)=1$, multiply the kernel series by the series
for the shorter word and integrate each term. For $p\ne-1$,
\begin{equation}
 \int X^p\ln^k X\,\dd X
 =X^{p+1}\sum_{j=0}^{k}
       \frac{(-1)^j k!}{(k-j)!(p+1)^{j+1}}\ln^{k-j}X.
\end{equation}
For $p=-1$ the primitive is $\ln^{k+1}X/(k+1)$.
The zero-letter words are normalized as
$G(\underbrace{0,\ldots,0}_{k};X)=\ln^k X/k!$.
All other integration constants follow from their vanishing limit at
the lower endpoint. This gives the required half-integer power series
without a numerical fit.

The rational coefficients of the total have poles of order at most
four at $X=0$. Accordingly the GPL series are computed through $X^8$
before multiplication, which suffices for the total through $X^4$.
On summing all terms, the negative powers, the constant and linear
terms, every half-integer power, and powers of $\ln X$ above one
cancel exactly. The first surviving term is the $X^2$ term displayed
in \eqref{eq:comp-collected-total-low-four}. The local kernel series
converge in a sufficiently small interval $0<X<1$.
After the displayed terms are removed, the next possible contributions
have the form $X^5$ times an analytic function or $X^5\ln X$ times an
analytic function, giving the stated remainder. Indeed, every
non-elementary word in the product form has at most one trailing zero,
and the only explicit factor $\ln X$ multiplying a GPL multiplies a
word with no trailing zero. The elementary part is also at most linear
in $\ln X$. Thus no higher logarithmic power can first appear in the
omitted orders.

Only after deriving these coefficients do we compare the leading
term with Lee, Schwartz and Zhang~\cite{qed-processes-lsz},
Equation (10). Their variable is the same $X=(s_{\mathrm{phys}}-m^2)/m^2$,
and their order-$\alpha^3$ normalization is $\alpha^3/m^2$.
The logarithmic coefficient $-16/9$ and the finite coefficient $7/15$
agree separately. That displayed literature equation does not give
the $X^3$ and $X^4$ coefficients, so no comparison of those two orders
is inferred from it.

